\newpage
\section{Giới thiệu}

Trong những năm gần đây, công nghệ Thực tế tăng cường (AR) đã trở thành một trong những hướng phát triển quan trọng trong lĩnh vực tương tác người -- máy (HCI), với ứng dụng ngày càng mở rộng trong giáo dục, y tế, thương mại điện tử, giải trí và đào tạo kỹ thuật.

Tuy nhiên, một trong những thách thức lớn nhất trong phát triển hệ thống AR là làm cho các đối tượng ảo xuất hiện tự nhiên và thuyết phục trong môi trường thực. Nếu các đối tượng ảo không tương tác hợp lý với không gian xung quanh -- chẳng hạn như không đứng đúng trên bề mặt thật, không bị che bởi các vật thể trong môi trường, hoặc có ánh sáng không phù hợp -- trải nghiệm AR sẽ trở nên thiếu chân thực và làm giảm hiệu quả của ứng dụng.

Để giải quyết vấn đề này, hệ thống AR hiện đại cần có khả năng \textit{Scene Understanding} và đảm bảo \textit{Realism} trong quá trình hiển thị nội dung ảo. Báo cáo này tập trung nghiên cứu bốn kỹ thuật cốt lõi phục vụ mục tiêu đó: ước lượng độ sâu và tái tạo lưới 3D của môi trường, xử lý che khuất giữa vật thể thật và vật thể ảo, ước lượng ánh sáng, và kết xuất bóng đổ.

Thông qua việc phân tích cơ chế hoạt động, phương pháp triển khai và thách thức của từng kỹ thuật, báo cáo hướng tới cung cấp một cái nhìn có hệ thống về vai trò của Scene Understanding trong việc xây dựng các hệ thống AR có khả năng hiển thị nội dung ảo một cách tự nhiên và thuyết phục trong môi trường thực.